MySQL to otwartoźródłowy relacyjny system zarządzania bazą danych rozwijany
od 1995 roku, początkowo przez szwedzką firmę MySQL AB, a obecnie będący
własnością Oracle Corporation. Obecna linia 8.0.x (stan na 2025~r.) oferuje
pełną obsługę transakcji ACID, silnik składowania InnoDB jako domyślny oraz
szeroką zgodność ze standardem SQL. MySQL jest dystrybuowany w dwóch wariantach
licencyjnych: GPL (Community Edition, bezpłatna) oraz komercyjnym (Enterprise
Edition, z dodatkowymi funkcjami bezpieczeństwa i narzędziami zarządzania).
Ze względu na dominującą pozycję w ekosystemie LAMP oraz domyślne wsparcie
ze strony platform takich jak WordPress czy Drupal, MySQL pozostaje jednym
z najszerzej stosowanych silników relacyjnych na świecie.

\subsection{Skalowalność}

PostgreSQL dobrze radzi sobie z rosnącą ilością danych, przede wszystkim przez
dodawanie zasobów do jednego serwera. Silnik potrafi efektywnie wykorzystać
dużą ilość RAM, wiele rdzeni procesora i szybkie dyski. Zapytania mogą być
wykonywane równolegle przez kilku pracowników jednocześnie, co szczególnie
przyspiesza operacje na dużych zbiorach danych. Dostępne jest też wbudowane
partycjonowanie tabel, czyli podział jednej dużej tabeli na mniejsze kawałki
według zakresu dat, wartości lub klucza, dzięki czemu baza nie przeszukuje
całej tabeli, tylko te partycje, które są potrzebne.

Skalowanie horyzontalne, czyli rozbudowa przez dodawanie kolejnych serwerów,
nie jest natywną cechą PostgreSQL. Można to osiągnąć przez replikację do
serwerów odczytu lub przez zewnętrzne rozszerzenia jak Citus, które dodają
automatyczny podział danych między węzły. Warto też pamiętać, że każde
połączenie z bazą to osobny proces systemowy, więc przy dużej liczbie
użytkowników zaleca się stosowanie pośredników jak PgBouncer.

\subsection{Bezpieczeństwo}

PostgreSQL oferuje rozbudowaną kontrolę dostępu. To kto i skąd może się
połączyć z bazą konfiguruje się w pliku \texttt{pg\_hba.conf}, gdzie definiuje
się kombinacje bazy, użytkownika i adresu IP oraz metodę uwierzytelniania.
Obsługiwane są hasła (domyślnie szyfrowane przez scram-sha-256), certyfikaty
klienckie, Kerberos, LDAP i inne. Połączenie może być szyfrowane przez SSL/TLS.

Uprawnienia do obiektów bazodanowych zarządza się przez role, które mogą
pełnić funkcję zarówno użytkowników jak i grup. Od wersji 9.5 dostępne jest
też Row-Level Security, które pozwala filtrować wiersze zwracane przez zapytanie
w zależności od tego, kto je wykonuje. Przydaje się to szczególnie w aplikacjach
obsługujących wielu klientów na jednej instancji bazy.

\subsection{Awaryjność}

Podstawą trwałości danych jest mechanizm Write-Ahead Log (WAL). Każda zmiana
trafia najpierw do dziennika, a dopiero potem do właściwych plików. Dzięki temu
w przypadku awarii serwera baza po restarcie automatycznie odtwarza stan sprzed
problemu i żadna zatwierdzona transakcja nie jest tracona.

Na bazie WAL działają mechanizmy replikacji i kopii zapasowych. Replikacja
strumieniowa pozwala utrzymywać aktualną kopię bazy na osobnym serwerze,
gotową do przejęcia ruchu w razie awarii. Point-in-Time Recovery umożliwia
odtworzenie bazy do dowolnego momentu w przeszłości. Automatyczny failover
realizuje się przez zewnętrzne narzędzia jak Patroni.

\subsection{Migracje}

Do przenoszenia danych między instancjami PostgreSQL służą \texttt{pg\_dump}
i \texttt{pg\_restore}, które tworzą przenośne zrzuty logiczne działające między
różnymi wersjami. Przy aktualizacji wersji na tym samym serwerze można użyć
\texttt{pg\_upgrade}, które nie przenosi danych, tylko aktualizuje struktury.
Migrację z innych silników, na przykład z MySQL, ułatwia narzędzie pgloader.

Wersjonowanie schematu bazy, czyli zarządzanie kolejnymi zmianami struktury
tabel, realizuje się przez zewnętrzne narzędzia jak Flyway lub Liquibase, albo
przez mechanizmy wbudowane w frameworki aplikacyjne. Warto podkreślić, że
PostgreSQL obsługuje transakcyjne DDL, co oznacza, że błędna migracja może
zostać wycofana tak jak zwykła transakcja.

\subsection{Integracje}

PostgreSQL posiada sterowniki dla praktycznie każdego popularnego języka
programowania i jest wspierany przez wszystkie główne frameworki ORM. Narzędzia
do analizy danych i wizualizacji jak Tableau, Grafana czy Metabase mają natywne
połączenia z PostgreSQL.

Ciekawą funkcją są Foreign Data Wrappers (FDW), które pozwalają odpytywać
zewnętrzne źródła danych, takie jak inne bazy czy pliki CSV, jakby były
zwykłymi tabelami. PostgreSQL jest też dostępny jako usługa zarządzana na
AWS, Google Cloud i Azure.


\subsection{Obszary biznesowe}

MySQL jest szczególnie mocno zakorzeniony w zastosowaniach webowych i jest
domyślnym wyborem dla wielu popularnych platform. Sprawdza się najlepiej
w następujących obszarach:

\begin{itemize}
    \item \textbf{Aplikacje webowe i CMS} -- WordPress, Drupal, Joomla oraz
    większość popularnych platform e-commerce jak Magento czy PrestaShop
    używa MySQL jako domyślnego backendu. Szacuje się, że MySQL napędza
    ponad 60\% stron internetowych korzystających z relacyjnych baz danych.
    \item \textbf{E-commerce} -- prostota konfiguracji, szeroka dostępność
    hostingu i dojrzałość ekosystemu sprawiają, że MySQL jest pierwszym
    wyborem dla sklepów internetowych małej i średniej skali.
    \item \textbf{Systemy SaaS} -- dzięki modelowi wątkowemu i dobrej
    wydajności przy dużej liczbie równoczesnych połączeń MySQL sprawdza się
    w aplikacjach obsługujących wielu użytkowników jednocześnie.
    \item \textbf{Startupy i projekty webowe} -- niska bariera wejścia,
    ogromna baza poradników i gotowych rozwiązań oraz dostępność na
    praktycznie każdym hostingu współdzielonym sprawiają, że MySQL
    jest naturalnym pierwszym wyborem dla nowych projektów.
    \item \textbf{Systemy analityki webowej} -- MySQL bywa stosowany
    do przechowywania i analizy logów, statystyk odwiedzin i danych
    behawioralnych użytkowników w zastosowaniach, gdzie nie jest wymagana
    zaawansowana analityka.
\end{itemize}

\subsection{Zalety}

\begin{itemize}
    \item \textbf{Prostota i niska bariera wejścia} -- MySQL jest łatwy
    w instalacji i konfiguracji. Dostępny jest na praktycznie każdym
    hostingu współdzielonym, co czyni go najbardziej dostępnym silnikiem
    relacyjnym dla początkujących.
    \item \textbf{Bardzo duży ekosystem} -- ogromna baza dokumentacji,
    poradników, gotowych narzędzi i społeczności sprawia, że znalezienie
    pomocy przy problemach jest łatwe.
    \item \textbf{Model wątkowy} -- w przeciwieństwie do PostgreSQL każde
    połączenie obsługiwane jest przez lekki wątek zamiast osobnego procesu,
    co zmniejsza narzut pamięciowy przy dużej liczbie połączeń.
    \item \textbf{Wysoka wydajność przy prostych zapytaniach} -- MySQL
    jest zoptymalizowany pod kątem szybkich odczytów i prostych operacji
    CRUD, co sprawdza się w typowych zastosowaniach webowych.
    \item \textbf{MySQL InnoDB Cluster} -- wbudowane rozwiązanie wysokiej
    dostępności z automatycznym failoverem bez potrzeby zewnętrznych narzędzi.
    \item \textbf{Szeroka dostępność w chmurze} -- dostępny jako usługa
    zarządzana na AWS, Google Cloud i Azure, a także w wariancie Aurora
    oferującym znacznie wyższą wydajność przy zachowaniu kompatybilności.
    \item \textbf{Kompatybilność} -- MySQL jest kompatybilny z wieloma
    narzędziami i frameworkami, a jego dialekt SQL jest szeroko znany
    wśród programistów.
\end{itemize}

\subsection{Wady i ograniczenia}

\begin{itemize}
    \item \textbf{Słabsze wsparcie standardu SQL} -- MySQL historycznie
    odchodzi od standardu SQL w wielu miejscach, co może powodować
    problemy przy migracji zapytań między silnikami. Przykładem jest
    domyślny tryb \texttt{ONLY\_FULL\_GROUP\_BY} wyłączony w starszych
    wersjach, który pozwalał na niepoprawne zapytania GROUP BY.
    \item \textbf{Brak transakcyjnego DDL} -- instrukcje \texttt{ALTER TABLE}
    czy \texttt{CREATE TABLE} niejawnie zatwierdzają bieżącą transakcję,
    co oznacza, że błędna migracja schematu nie może zostać wycofana
    przez \texttt{ROLLBACK}.
    \item \textbf{Ograniczona analityka} -- MySQL nie obsługuje wielu
    zaawansowanych funkcji analitycznych SQL takich jak pełne wsparcie
    dla funkcji okienkowych (do niedawna), \texttt{GROUPING SETS},
    \texttt{CUBE} czy \texttt{ROLLUP}. PostgreSQL jest znacznie silniejszy
    w zastosowaniach HTAP.
    \item \textbf{Brak Row-Level Security} -- filtrowanie wierszy per
    użytkownik musi być implementowane na poziomie aplikacji lub widoków,
    co zwiększa złożoność systemów wielodostępnych.
    \item \textbf{Słabsze typy danych} -- MySQL oferuje uboższy katalog
    typów w porównaniu do PostgreSQL. Brakuje natywnych typów zakresowych,
    tablic czy zaawansowanych typów geometrycznych bez zewnętrznych rozszerzeń.
    \item \textbf{Właściciel korporacyjny} -- MySQL należy do Oracle,
    co dla części organizacji jest czynnikiem ryzyka przy wyborze silnika
    dla krytycznych systemów. Część zaawansowanych funkcji dostępna jest
    wyłącznie w płatnej Enterprise Edition.
    \item \textbf{Ograniczone indeksowanie} -- MySQL obsługuje B-tree,
    Hash i FULLTEXT, ale brakuje zaawansowanych typów indeksów dostępnych
    w PostgreSQL (GIN, GiST, BRIN), co ogranicza możliwości przy
    wyszukiwaniu pełnotekstowym i danych przestrzennych.
\end{itemize}

\subsection{Udogodnienia}

\begin{itemize}
    \item \textbf{MySQL Workbench} -- oficjalne graficzne narzędzie do
    administracji, projektowania schematów (wizualny edytor ERD), migracji
    i profilowania zapytań. Dostępne na Windows, macOS i Linux.
    \item \textbf{Narzędzia CLI} -- \texttt{mysql} (interaktywny klient),
    \texttt{mysqldump} (eksport), \texttt{mysqlcheck} (weryfikacja tabel),
    \texttt{mysqladmin} (administracja serwerem).
    \item \textbf{EXPLAIN i query profiler} -- MySQL oferuje \texttt{EXPLAIN}
    i \texttt{EXPLAIN ANALYZE} (od wersji 8.0) do analizy planów wykonania,
    a także \texttt{performance\_schema} i \texttt{sys} schema do
    monitorowania wydajności na poziomie zapytań, wątków i operacji I/O.
    \item \textbf{JSON} -- od wersji 5.7 MySQL obsługuje natywny typ JSON
    z operatorami ścieżkowymi (\texttt{->}, \texttt{->>}) i możliwością
    tworzenia indeksów na wygenerowanych kolumnach opartych o pola JSON.
    \item \textbf{Generowane kolumny} -- kolumny wyliczane automatycznie
    na podstawie wyrażenia (wirtualne lub przechowywane), przydatne do
    indeksowania fragmentów JSON lub wyrażeń obliczeniowych.
    \item \textbf{Szeroka dostępność hostingu} -- MySQL jest dostępny
    na praktycznie każdym hostingu współdzielonym i VPS, co czyni go
    najbardziej dostępnym silnikiem relacyjnym pod względem infrastruktury.
\end{itemize}